\documentclass[12pt,a4paper]{article}

\usepackage[utf8]{inputenc}
\usepackage[T5]{fontenc}
\usepackage[utf8]{vntex}

\usepackage{lmodern}
\usepackage{geometry}
\geometry{margin=2.5cm}
\usepackage{graphicx}
\graphicspath{{figures/}{./}}
\usepackage{amsmath}
\usepackage{booktabs}
\usepackage{float}
\usepackage{array}
\setlength{\parindent}{0pt}
\setlength{\parskip}{6pt}
\usepackage{caption}
\usepackage{subcaption}
\usepackage{xcolor}
\usepackage[hypertexnames=false]{hyperref}
\usepackage{longtable}

% English names. vntex installs Vietnamese caption names from a
% begin-document hook, so ours must be deferred to the same hook to win.
\AtBeginDocument{%
  \renewcommand{\contentsname}{Contents}%
  \renewcommand{\listfigurename}{List of Figures}%
  \renewcommand{\listtablename}{List of Tables}%
  \renewcommand{\figurename}{Figure}%
  \renewcommand{\tablename}{Table}%
  \renewcommand{\refname}{References}%
}

% Keep floats inside their section
\usepackage[section]{placeins}

\hypersetup{
    colorlinks=true,
    linkcolor=black,
    urlcolor=blue,
}

\setlength{\floatsep}{10pt plus 2pt minus 2pt}
\setlength{\textfloatsep}{12pt plus 2pt minus 4pt}
\setlength{\intextsep}{10pt plus 2pt minus 2pt}
\renewcommand{\textfraction}{0.05}
\renewcommand{\floatpagefraction}{0.85}
\setlength{\emergencystretch}{3em}
\hbadness=2000

% Figure helper: insert the real image when the file is present, otherwise a
% clean placeholder box, so the report compiles before the diagrams exist.
% Drop the named .png into this folder later and the figure upgrades itself.
\newcommand{\reportfig}[3]{%
  \IfFileExists{figures/#2}{%
    \includegraphics[width=#1\textwidth]{#2}%
  }{%
    \fbox{\begin{minipage}[c][5.5cm][c]{#1\textwidth}%
      \centering\itshape
      [\,Figure to be inserted\,]\par\vspace{6pt}
      #3\par\vspace{10pt}%
      \normalfont\ttfamily\footnotesize #2%
    \end{minipage}}%
  }%
}

% Logo helper: same idea for the title-page logos, so a missing logo file does
% not break compilation.
\newcommand{\titlelogo}[1]{%
  \IfFileExists{#1}{\includegraphics[width=0.9\textwidth]{#1}}{%
    \fbox{\parbox[c][2.6cm][c]{0.8\textwidth}{\centering\ttfamily\footnotesize #1}}%
  }%
}

\begin{document}
\raggedbottom
\pagenumbering{gobble}

% =========================================================================
% TITLE PAGE
% =========================================================================
\begin{titlepage}
    \centering
    \textbf{\large University of Engineering and Technology, VNU, Hanoi}\\[0.2cm]
    \textbf{\large Institute for Artificial Intelligence}\\[0.5cm]
    \rule{10cm}{0.4pt}\\

    \vspace{2cm}
    \begin{minipage}{0.45\textwidth}
        \begin{center}
            \titlelogo{uet_logo.jpg}
        \end{center}
    \end{minipage}
    \hfill
    \begin{minipage}{0.45\textwidth}
        \begin{center}
            \titlelogo{ai_logo.png}
        \end{center}
    \end{minipage}

    \vspace{1.5cm}
    {\large \textbf{COURSE REPORT}}\\[0.5cm]
    {\large \textbf{INDUSTRIAL INTERNSHIP}}\\[0.5cm]

    \vspace{5pt}
    {\large{
    Topic: \textbf{HumJob: A Local-First Browser Toolkit for the Singing
    Voice - Transcription, Analysis, and Ear/Vocal Training}
    }}\\[0.25cm]

    \vspace{12pt}
    {\large
    \begin{tabular}{>{\bfseries}l l}
        Supervisor: & PhD. Trần Hồng Việt \\
    \end{tabular}
    }\\[0.5cm]

    \vspace{12pt}
    {\large
    \begin{tabular}{>{\bfseries}l l}
        Author: & Vũ Nguyên Đan - 23020351 \\
    \end{tabular}
    }\\[0.5cm]

    \vfill
    \centering
    \textbf{Hà Nội, 2026}
\end{titlepage}

% =========================================================================
% EVALUATION PAGE
% =========================================================================
\newpage
\begin{center}
    \textbf{Supervisor's Evaluation:}
\end{center}

\vspace{0.2cm}

\noindent\makebox[\linewidth]{\dotfill}\\[0.2cm]
\noindent\makebox[\linewidth]{\dotfill}\\[0.2cm]
\noindent\makebox[\linewidth]{\dotfill}\\[0.2cm]
\noindent\makebox[\linewidth]{\dotfill}\\[0.2cm]
\noindent\makebox[\linewidth]{\dotfill}\\[0.2cm]
\noindent\makebox[\linewidth]{\dotfill}\\[0.2cm]
\noindent\makebox[\linewidth]{\dotfill}\\[0.2cm]
\noindent\makebox[\linewidth]{\dotfill}\\[0.2cm]
\noindent\makebox[\linewidth]{\dotfill}\\[0.2cm]
\noindent\makebox[\linewidth]{\dotfill}\\[0.2cm]
\noindent\makebox[\linewidth]{\dotfill}\\[0.2cm]
\noindent\makebox[\linewidth]{\dotfill}\\[0.2cm]
\noindent\makebox[\linewidth]{\dotfill}\\[0.2cm]
\noindent\makebox[\linewidth]{\dotfill}\\[0.2cm]
\noindent\makebox[\linewidth]{\dotfill}\\[0.2cm]
\noindent\makebox[\linewidth]{\dotfill}\\[0.2cm]
\noindent\makebox[\linewidth]{\dotfill}\\[0.2cm]
\noindent\makebox[\linewidth]{\dotfill}

\noindent\hspace{3cm} Score (in figures) \dotfill{} Score (in words) \dotfill

\vspace{1.5cm}

\begin{center}
    Hà Nội, \ldots{} \ldots{} 2026\\[0.2cm]
    \textbf{Supervisor}\\[0.1cm]
    (Signature and full name)
\end{center}

% =========================================================================
% ACKNOWLEDGEMENTS
% =========================================================================
\newpage
\begin{center}
    \textbf{\large ACKNOWLEDGEMENTS}
\end{center}

I would like to thank my supervisor for the guidance and feedback received
throughout the Industrial Internship course, which shaped both the direction of this work,
\textbf{HumJob}, and the standard of intellectual honesty I have tried to hold
it to.

This report is the work of a single student. Given the breadth of the system
and the limited time available, it inevitably contains shortcomings. In
particular, the central finding of the transcription evaluation is a sobering
one: the aggregate synthetic benchmark the transcriber was tuned against
saturated at ceiling, while the system is not yet validated on real hummed input.
I have chosen to report that honestly, and to rebuild the evaluation around richer
evidence, rather than to present a saturated benchmark as a success. I welcome any
comments and corrections from the instructor so that the work may be improved.

% =========================================================================
% ABSTRACT
% =========================================================================
\newpage
\begin{center}
    \textbf{\large ABSTRACT}
\end{center}

\textbf{HumJob} is a local-first, browser-based toolkit for the singing voice.
It brings together six tools that share one foundation, pitch detection applied
to a human voice, and that all run on the user's own machine: a \emph{melody
transcriber} that turns a hummed tune into notes, key, and suggested chords and
exports them as MIDI, MusicXML, and engraved sheet music; a \emph{pitch finder}
that reports the key, tempo, and Camelot code of any clip, including polyphonic
music; a \emph{transposer} that shifts a score to a new key; a \emph{real-time}
pitch monitor and vocal trainer; a \emph{sing-along} trainer that scores a sung
performance against a reference melody; and an \emph{ear trainer} for interval,
chord, scale, and scale-degree recognition. A \emph{practice hub} folds the
training history the tools record locally into progress trends. Several features
add optional, opt-in language-model assistance, spoken-language coaching, a
weekly practice plan, natural-language score edits, a whole-melody rewrite, and
reharmonisation, which is the only part of the system that ever leaves the user's
machine.

The technical core of the toolkit, and the subject of this report's formal
evaluation, is the transcriber. It solves monophonic melody transcription from a
bare human voice, which is hard because a continuous pitch contour has to be cut
into discrete notes with correct pitches, onsets, and durations. HumJob makes
this manageable with two deliberate user-interface constraints. The user hums the
syllable ``da-da-da'', whose consonant gives every note a crisp onset and
separates repeated notes of the same pitch; and the user hums to a metronome at
a known tempo, which turns rhythm from estimation into snapping onto a known
grid. The pipeline is built as a chain of pure functions over three data types,
with every tunable parameter centralised, and supports five interchangeable
pitch backends (pyin, CREPE, PESTO, FCNF0++, and basic-pitch). A grid-aware
segmenter and a consolidation stage that runs for every backend address the dominant
failure mode, and a spectral octave-correction stage repairs a systematic
pitch-tracker error on legato lines.

The transcriber is evaluated on a synthetic fixture harness that scores note
detection (pitch and onset F1) and, separately, rhythm (quantised onset and
duration accuracy against an intended grid). On the expressive ``realistic''
fixtures the note F1 rose from $0.799$ to $1.000$ across three engineering
iterations, and the harness saturated; the rhythm half saturated too once a
tempo-refinement stage was added. A single F1 at ceiling cannot separate a better
segmenter from a worse one, and it is entirely synthetic, so the evaluation was
rebuilt around four kinds of evidence that need no precise singer, since the
author is not one and no labelled real-hum dataset exists. These are: diagnostic
error-rates over a larger, diversified corpus, which discriminate again and
localise the dominant expressive-input failure to merged same-pitch repeats
(repeat recall falls to about $0.61$); metamorphic correctness properties
(transpose, gain, noise, silence-pad, and tempo invariance) that need no reference
melody and hold on real audio too; a reconstruction round-trip that scores a real
hum against its own resynthesised audio (an unsupervised sanity signal, mean
$0.62$ over four real takes); and a calibration of the synthetic realism against
those real takes, which confirms the assumed vibrato rate and drift are close
to real voice. A useful, honest finding is that this pipeline does not
over-split under heavy perturbation, as had been expected; its defences hold and
it instead degrades by merging or dropping notes. None of the four is a
real-world accuracy number, which still requires real ground truth, but together
they replace one saturated number with a picture that stays useful. This report
documents the whole toolkit, concentrates its evaluation on the transcriber, and
gives an honest account of what each kind of evidence can and cannot demonstrate.

% =========================================================================
% TABLE OF CONTENTS
% =========================================================================
\newpage
\pagenumbering{arabic}
\tableofcontents
\newpage
\listoffigures
\newpage
\listoftables
\newpage

% =========================================================================
% SECTION 1: INTRODUCTION
% =========================================================================
\section{Introduction}

\subsection{Motivation and significance}

A common musical situation has no good tool. A person has a melody in their
head, hums it, and wants it written down: as notes on a staff, as a MIDI file
they can edit, or as a chord sketch they can play. Turning that hum into
notation is the task known as \emph{query-by-humming} transcription, and it is
harder than it looks. The human voice produces a single continuous pitch
contour that slides, wavers, and drifts; notation requires that contour to be
cut into a sequence of discrete notes, each with a definite pitch, a definite
start time, and a definite written duration, and it requires a key and a
metrical grid within which those durations are meaningful.

Two families of existing tools sidestep parts of the problem rather than
solving it. General audio-to-MIDI transcribers are trained on instruments and
polyphonic music; on a bare voice they are prone to octave jumps and they do
not exploit any of the structure a deliberate hum provides. Real-time pitch
monitors and tuners report the instantaneous pitch but never commit to a note
segmentation at all. What is missing is a system that is built specifically for
a cooperative human who is willing to hum in a way that makes transcription
manageable.

HumJob is that system. Its central design decision is to accept two small
constraints on how the user hums, and to build the entire pipeline around them:

\begin{itemize}
    \item The user hums the syllable \textbf{``da-da-da''}. The consonant
    closes the vocal tract briefly at the start of each note, which produces a
    crisp amplitude onset and, crucially, separates two repeated notes of the
    same pitch that a legato hum would smear into one.
    \item The user hums \textbf{to a metronome at a known tempo}. A known
    tempo grid converts rhythm from a hard estimation problem into a snapping
    problem: note onsets and durations are aligned to the nearest grid
    position rather than guessed from raw timing.
\end{itemize}

These constraints are load-bearing. They are what make the problem manageable
for a single-student project, and a recurring theme of this report is that the
remaining failures of the transcriber are failures to exploit them fully, not
failures of the underlying pitch models.

\subsection{From a transcriber to a toolkit for the voice}
\label{sec:toolkit-intro}

HumJob began as that transcriber, which remains its technical core and the focus
of this report's evaluation. In the course of building it, a family of closely
related tools for the singing voice fell out of the same foundation, because
every one of them rests on the same operation: estimating the pitch of a human
voice, locally and in the browser. The finished system is therefore a toolkit of
six tools, presented as tabs in one local-first web application:

\begin{itemize}
    \item a \textbf{Transcriber} that turns a hum into notes, key, and chords,
    the core pipeline, with an interactive staff editor for correcting its draft
    and offline or language-model reharmonisation of its chords;
    \item a \textbf{Pitch Finder} that reports the key, tempo, and Camelot code
    of an uploaded clip, including polyphonic music;
    \item a \textbf{Transposer} that shifts a score, from a file or from the last
    hum, to a new key;
    \item a \textbf{Realtime} monitor and vocal trainer: a live pitch display, a
    guitar tuner, and a set of practice drills;
    \item a \textbf{Sing-Along} trainer that plays a reference melody and scores
    a sung performance note by note, with optional language-model coaching;
    \item an \textbf{Ear Trainer} that plays intervals, chords, scales, and
    cadences for the user to identify.
\end{itemize}

A seventh tab, a \textbf{practice hub}, is a home dashboard rather than a tool in
this sense: it aggregates the training history the other tabs record locally into
offline progress trends. Several tabs also expose optional, opt-in
language-model assistance, described in Section~\ref{sec:coaching}.

This report describes the whole toolkit but concentrates its evaluation on the
transcriber, because that is where the project's research contribution and its
hardest, most measurable problem lie. The other five tools are engineering
deliverables built on the shared foundation; Section~\ref{sec:toolkit} describes
them and Appendix~A demonstrates each.

\subsection{Objectives}

The project pursues the following objectives:

\begin{itemize}
    \item To design and implement a complete, local-first melody-transcription
    pipeline that takes a hummed recording and returns notes, key, and
    suggested chords, exported as MIDI, MusicXML, and engraved sheet music.
    \item To structure the pipeline as a chain of pure functions over a small,
    fixed data model, so that every stage is independently testable and every
    tunable parameter is centralised rather than scattered as magic numbers.
    \item To support several interchangeable pitch-detection backends behind a
    single interface, so the note-production model can be chosen or replaced
    without touching the rest of the pipeline.
    \item To deliver the pipeline both as a command-line tool and as a browser
    web application that records audio locally and calls the same code, with an
    interactive editor for correcting the automatic draft.
    \item To reuse the same local-first pitch-detection foundation across a
    family of related tools for the singing voice, analysis, transposition,
    real-time monitoring, and vocal and ear training, so that the core serves
    more than one musical task.
    \item To build an evaluation harness that measures the transcriber honestly,
    and to characterise clearly what that harness can and cannot demonstrate.
\end{itemize}

\subsection{Research questions}

The project is organised around three research questions. They concern the
transcriber specifically, since that is where the investigative work lies; the
other tools are assessed as functional deliverables
(Section~\ref{sec:toolkit}, Appendix~A).

\begin{itemize}
    \item \textbf{RQ1 - Do the interface constraints help?} Do the two
    deliberate constraints (a consonant onset per note, and humming to a known
    tempo grid) turn melody transcription from an estimation problem into a
    manageable segmentation and snapping problem?
    \item \textbf{RQ2 - Is this a segmentation problem or a model problem?}
    When notes come out wrong, does the fault lie chiefly in the segmentation
    and quantisation of an otherwise-correct pitch contour, rather than in the
    pitch model itself?
    \item \textbf{RQ3 - How well can the system be evaluated?} How well does
    the system agree with ground truth, and what are the methodological limits
    of that measurement given that the project has only synthetic ground truth
    and no labelled corpus of real hummed melodies?
\end{itemize}

% =========================================================================
% SECTION 2: BACKGROUND
% =========================================================================
\section{Background}

\subsection{Melody transcription and the monophonic subproblem}

Automatic music transcription in general is the conversion of an audio signal
into a symbolic score. The general problem is polyphonic and very hard. HumJob
addresses the restricted \emph{monophonic} case: one voice, one note at a time.
This restriction removes the need to separate simultaneous pitches, but it does
not make the problem trivial, because the difficulty moves from
source-separation to \emph{segmentation}: deciding where one note ends and the
next begins, and how long each note is in metrical terms.

A monophonic transcription pipeline has three logical stages: estimate a pitch
contour over time, cut that contour into discrete note events, and place those
events on a key and a rhythmic grid. HumJob follows exactly this shape, and the
bulk of its engineering effort is in the second and third stages.

\subsection{Fundamental-frequency estimation}

The first stage estimates the fundamental frequency ($f_0$) of the voice, frame
by frame. Several families of estimator exist and HumJob supports five of them
behind one interface:

\begin{itemize}
    \item \textbf{pYIN}~\cite{pyin} is a classical signal-processing estimator
    that augments the YIN algorithm with probabilistic thresholds and a Viterbi
    decoding step. It is fast, dependency-light, and the project default.
    \item \textbf{CREPE}~\cite{crepe} is a convolutional neural network trained
    to estimate pitch directly from the waveform, and is steady on singing
    voice.
    \item \textbf{PESTO}~\cite{pesto} is a self-supervised,
    transposition-equivariant pitch estimator; it is the web application's
    default because it is precise on voice while being lighter and faster than
    CREPE.
    \item \textbf{FCNF0++}~\cite{penn} is a fully-convolutional pitch and
    periodicity estimator, a precision peer of PESTO.
    \item \textbf{basic-pitch}~\cite{basicpitch} is an instrument-trained
    neural note-transcription model that emits note events directly, bypassing
    the frame-level contour.
\end{itemize}

The important background fact for this report concerns a specific failure of
Viterbi-decoded trackers. On a continuously voiced legato line, with no silence
to reset the decoder, the ``stay put'' transition prior of a Viterbi decoder
can lock a whole note onto its subharmonic, reporting it a full octave too low.
This is not an audio artifact; it is a decoding artifact, and Section~%
\ref{sec:octave} describes how HumJob detects and repairs it.

\subsection{From a pitch contour to notes}

Turning a contour into notes requires three sub-decisions. \emph{Voicing}
decides which frames carry pitch at all, gating out breaths and silence.
\emph{Segmentation} decides where the boundaries between notes fall, which for
a ``da-da-da'' hum means finding the brief energy dips created by the
consonant. \emph{Consolidation} does the opposite where needed: a single held
note that vibrato has fragmented into several pieces must be fused back into
one. Segmentation and consolidation pull in opposite directions, and balancing
them is the core difficulty of the pipeline.

\subsection{Key finding and rhythm quantisation}

Once notes exist, two musical decisions remain. \emph{Key finding} identifies
the scale the melody sits in; HumJob correlates the sung pitch-class
distribution against the Krumhansl-Kessler probe-tone profiles for all
twenty-four major and minor keys~\cite{krumhansl}. \emph{Quantisation} snaps
note onsets and durations to a metrical grid derived from the known tempo,
which is what allows durations to be written as clean note values rather than
as awkward fractions. When the tempo is wrong, quantisation produces
non-integer durations that notation software renders as strings of tied notes,
which is the visible symptom of most rhythm errors.

\subsection{Why a metronome and ``da-da-da''}

The two constraints introduced in Section~1 are the project's answer to the two
hardest sub-problems above. The consonant in ``da-da-da'' directly supplies the
onset cue that segmentation needs, and in particular it is the only reliable
way to separate two consecutive notes of the same pitch. The metronome supplies
the grid that quantisation snaps to. Neither constraint is a burden for a
cooperative user, and together they replace two estimation problems with two
detection problems. The project treats them as fixed assumptions, not as
limitations to be optimised away.

% =========================================================================
% SECTION 3: SYSTEM DESIGN AND METHODOLOGY
% =========================================================================
\section{System Design and Methodology}

\subsection{Architecture overview}

The transcriber is organised as a chain of pure functions over three data types.
Figure~\ref{fig:arch} shows the high-level structure: a core Python package
implements the pipeline, a command-line tool calls it directly, and a FastAPI
web application records audio in the browser and calls the same pipeline
server-side. The same web application also hosts the other five tools of the
toolkit as separate browser tabs (Section~\ref{sec:toolkit}); this section
describes the transcription pipeline, which is the shared technical core.
Everything runs locally, and no audio is ever sent to any third-party service. The exception to local-only operation is a small family of opt-in
AI features (Section~\ref{sec:coaching}): only when the user explicitly presses a
button does one of them send a short text or numeric summary, never any audio, to
an external language-model API.

\begin{figure}[H]
\centering
\reportfig{0.9}{architecture.png}{System architecture: the core pipeline
package, the command-line entry point, and the FastAPI web application with its
static browser front-end. The same pipeline code serves both entry points.}
\caption{High-level system architecture of HumJob.}
\label{fig:arch}
\end{figure}

The three data types are the contract that every stage shares
(Table~\ref{tab:types}). A \texttt{Frame} is one dense analysis frame, roughly
one per $11.6$ ms hop, carrying the frame time, the estimated $f_0$, a voicing
confidence, and a short-time energy. A \texttt{NoteEvent} is one discrete note
after segmentation, carrying its start and end in seconds, its integer pitch,
its raw fractional pitch, a cents offset from the nearest semitone, and, once
quantised, its grid-aligned position and duration in quarter-note units. A
\texttt{Score} is the whole transcription, ready to export.

\begin{table}[H]
\centering
\caption{The three data types the pipeline is built on.}
\label{tab:types}
\begin{tabular}{l p{0.60\textwidth}}
\toprule
\textbf{Type} & \textbf{Role} \\
\midrule
\texttt{Frame}     & One analysis frame ($\approx 11.6$ ms): frame time, $f_0$
(NaN if unvoiced), voicing confidence, short-time energy. What the pitch
tracker sees. \\
\texttt{NoteEvent} & One discrete note after segmentation: start/end seconds,
integer MIDI pitch, raw fractional pitch, cents offset, and (after
quantisation) grid position and duration in quarter-note units. \\
\texttt{Score}     & The whole transcription: the note sequence, the detected
key, tuning offset, and chords, ready to export. \\
\bottomrule
\end{tabular}
\end{table}

A design rule the project holds to is that every tunable parameter lives in a
single central parameter object, not as a magic number inside a stage. Changing
behaviour means changing a named knob, which keeps the stages readable and the
experiments reproducible.

\subsection{The transcription pipeline}

Figure~\ref{fig:pipeline} shows the pipeline stages in order. The chain is:

\begin{center}
audio $\rightarrow$ preprocess $\rightarrow$ [note production] $\rightarrow$
consolidate $\rightarrow$ tuning $\rightarrow$ key $\rightarrow$ quantise
$\rightarrow$ chords $\rightarrow$ export.
\end{center}

\begin{figure}[H]
\centering
\reportfig{0.92}{pipeline-transcriber.png}{The transcription pipeline. Note
production has two paths: the per-frame trackers (pyin / CREPE / PESTO /
FCNF0++) feed octave correction, voicing, and the grid-aware segmenter, while
basic-pitch emits note events directly. Both paths then share consolidation,
tuning, key finding, quantisation, chords, and export.}
\caption{The HumJob transcription pipeline, both note-production paths.}
\label{fig:pipeline}
\end{figure}

\textbf{Note production} has two paths, selected by a backend parameter. The
neural note-event backend (basic-pitch) emits notes directly and has no frame
contour. The per-frame trackers (pyin, CREPE, PESTO, FCNF0++) produce a
\texttt{Frame} contour that then passes through octave correction, a voicing
gate, and the grid-aware segmenter to become notes.

\textbf{Consolidation} runs for every backend. It fuses the fragments that a
held, vibrato-laden note leaves behind, which is the fix for the ``one note
becomes many slivers'' failure. It is grid-aware: it will not fuse two
same-pitch notes across a grid onset, so it never undoes a re-articulation that
the segmenter deliberately split on the beat.

\textbf{The remaining stages} apply a global tuning correction, find the key by
the Krumhansl-Kessler correlation, quantise onsets and durations onto the
tempo grid, suggest chords, and export to MIDI, MusicXML, and engraved sheet
music.

\subsection{The note-detection backends}

The five backends trade precision, speed, and training domain differently.
Table~\ref{tab:backends} summarises them. The default differs by entry point,
which is intentional: the library default is pyin (fast and dependency-light),
the command-line default is basic-pitch, and the web application default is
PESTO (most precise on voice). The four per-frame trackers all feed the same
segmenter; only basic-pitch bypasses it.

\begin{table}[H]
\centering
\caption{The note-detection backends.}
\label{tab:backends}
\begin{tabular}{l p{0.34\textwidth} p{0.34\textwidth}}
\toprule
\textbf{Backend} & \textbf{What it is} & \textbf{Best for / notes} \\
\midrule
\texttt{pesto} & Self-supervised, transposition-equivariant pitch estimator &
Most precise on voice; web-app default; lighter and faster than CREPE \\
\texttt{fcnf0} & FCNF0++, a fully-convolutional $f_0$ estimator & Precision
peer of PESTO; shines on real voice; optional install \\
\texttt{crepe} & Neural pitch CNN, voice-trained & Steady on humming and
singing; uses the accurate full model with Viterbi decoding \\
\texttt{basic\_pitch} & Instrument-trained note-transcription CNN & Instrument
clips; can octave-jump on a bare voice; bypasses the segmenter \\
\texttt{pyin} & Classical signal-processing $f_0$ plus the project segmenter &
Crisp staccato ``da-da-da''; library default \\
\bottomrule
\end{tabular}
\end{table}

\subsection{Grid-awareness: segmentation and consolidation}

The single most important methodological idea in the pipeline is that
segmentation is \emph{grid-aware}. The segmenter is passed the known tempo and
a fine short-window energy envelope. A short or soft consonant closure between
two same-pitch notes appears as a narrow dip in that fine envelope; a width
gate separates such a closure from a wide tremolo swell, and a shallower dip is
still accepted when it lands on a beat, because the grid says a note boundary is
expected there. This is what separates repeats of the same pitch, and it
replaced an earlier coarse valley splitter that smeared over short closures.

Consolidation is the necessary balancing force. Because vibrato and breath can
fragment one held note into several pieces, a consolidation
stage that runs for every backend fuses same-pitch fragments back together. Its grid-awareness is the safeguard: it refuses to fuse across a grid onset, so segmentation and
consolidation cannot fight each other into an oscillation.

\subsection{Spectral octave correction}
\label{sec:octave}

A distinct stage, inserted between the tracker and the voicing gate for the
per-frame backends, repairs the Viterbi subharmonic error described in
Section~2. The stage runs for every backend and rests on a spectral signature. A
genuine fundamental at frequency $f$ has energy at both its odd harmonics
($f, 3f, 5f$) and its even harmonics ($2f, 4f, 6f$). A subharmonic error, in
which the tracker has reported half the true frequency, has energy only at the
even multiples, because those are the true fundamental and its harmonics; the
odd multiples of the reported frequency have little or no energy. The stage
therefore computes, per voiced frame, the salience at the odd and even
multiples of the reported $f_0$, and doubles the reported frequency only when
the odd-harmonic salience collapses relative to the even-harmonic salience. The
test is deliberately conservative: it leaves genuine fundamentals alone,
including missing-fundamental voices, which retain energy at their odd
harmonics. Section~\ref{sec:results-note} reports the effect of this stage on
the evaluation.

\subsection{The interactive editor (Manual mode)}

The automatic draft is not always correct, and the weak link is segmentation
rather than the pitch contour. The web application therefore provides a Manual
mode: an in-browser, staff-based editor for fixing the automatic draft. It is
almost entirely client-side. It re-implements the score-building and
enharmonic-spelling logic of the export stage in the browser, engraves with a
vendored WebAssembly build of the Verovio music-notation
engine~\cite{verovio}, and offers pure editing operations (change pitch, change
duration, merge, split, delete to rest, insert) with undo and redo. A reference
pitch strip shows the hummed contour against the chosen notes, so the user can
see where the draft diverged from what they sang. The client sends the edited
notes back to the server only at the edges, to re-export MIDI and MusicXML and
to recompute key and chords.

Manual mode also accepts plain-language edits. An ``Ask'' box takes an
instruction such as ``merge notes 3 to 5'' or ``raise the last two notes an
octave''; the common phrasings are matched by an offline grammar that runs with
no network access, and anything else is turned into a small edit script by an
opt-in language model (Section~\ref{sec:coaching}). Either way the instruction
becomes the same validated script, which a pure interpreter in the browser
applies as one undoable step, rejecting the whole script if any note reference is
out of range. A separate ``Reorganise'' button goes the other way: rather than a
bounded local edit, it rewrites the whole melody into clean, easy-to-read note
values in one undoable step (Section~\ref{sec:coaching}).

The chords under the draft can be reworked from the result card, independently of
the note editing. Clicking a chord offers deterministic alternatives, computed
offline: for that measure the system ranks candidate chords, within a chosen
style set of triads, diatonic sevenths, and secondary dominants, by how much of
the measure's melody each one covers, with a small complexity penalty so a plain
triad wins a tie. A ``Reharmonise'' control instead asks the opt-in language
model for a whole new progression in a requested style (``jazzier'', ``simpler
for guitar''); the model proposes one chord per measure from a fixed quality
whitelist, and the system validates the proposal, spells and Roman-labels it with
music21~\cite{music21}, scores each bar's melody fit so a clash is shown rather
than hidden, and re-engraves the sheet. Applying either result updates the sheet,
the playback, the editor, and the transposer together.

\subsection{The wider toolkit: analysis, training, transposition, and ear
training}
\label{sec:toolkit}

The web application is tabbed. Besides the Transcriber described above, five tabs
are separate tools that do not go through the transcription segmenter; together
with the transcriber they make up the toolkit introduced in
Section~\ref{sec:toolkit-intro}. Each reuses the shared foundation, pitch
detection and the Web Audio and music-notation stack, for a different musical
task. The four voice-input tools below share one design rule with the
transcriber: their pure logic is separated from the DOM so it can be unit-tested,
and the microphone is released whenever a tab is left.

\begin{itemize}
    \item \textbf{Pitch Finder} analyses an uploaded clip for key, tempo, and a
    Camelot code, using a chroma-based method that works on polyphonic audio.
    It reuses the key-scoring and tempo-folding code but not the monophonic
    segmenter, which would be wrong for a full song.
    \item \textbf{Realtime} is a client-side live pitch monitor and guitar
    tuner built on the Web Audio API, with a set of vocal-training drills (a
    target-note lane, a match game, a scale trainer, vibrato analysis, a vocal
    range finder, and a per-session progress log). A server round-trip cannot
    be real-time, so this path runs entirely in the browser and only asks the
    server for the sung key on stop, from a pitch histogram, with no audio
    upload. The vocal range finder is the one drill that makes an empirical
    claim about the singer relative to a population: from the measured maximum
    phonational range it reports a percentile (\emph{wider than X\% of adults}).
    To keep that claim honest, the model's mean is taken from a cited study of
    adult phonational ranges~\cite{hollien}, while its standard deviation is a
    clearly-labelled assumption (about four semitones, in line with
    voice-range-profile literature) rather than a measured value, and both are
    printed in the interface.
    \item \textbf{Transposer} shifts a score to a new key. A file mode
    transposes an uploaded MIDI or MusicXML file server-side with
    music21~\cite{music21}, moving every voice and the key signature; a hum
    mode transposes the last transcription client-side. The file mode also
    offers Camelot-compatible key presets as one-click shifts.
    \item \textbf{Sing-Along} turns an uploaded MIDI or MusicXML score into a
    karaoke exercise: the melody is played back as a guide while the user's
    voice is tracked live and scored note by note. Because the whole melody is
    known in advance, the score itself drives the clock; there is no free-timing
    alignment. The server reduces a possibly polyphonic upload to a single
    monophonic melody line (its skyline, with ties stripped), and the browser
    schedules the count-in and the guide notes at absolute audio-clock times and
    scores the sung pitch against the active note. The in-tune tolerance is
    selectable (four difficulty levels), the guide volume is adjustable, and the
    take can be paused and resumed.
    \item \textbf{Ear Trainer} is the mirror of the other tabs: instead of the
    user making a sound for the system to analyse, the system plays a sound for
    the user to identify. It is fully client-side, with no microphone and no
    server. Four sub-modes, selected like the Realtime drills, cover interval,
    chord-quality, scale-or-mode, and scale-degree recognition; each question
    plays the same piano samples the other tabs use and offers multiple-choice
    answers, with three difficulty pools and a per-mode accuracy history kept in
    the browser. The key sub-mode tests scale-degree recognition after a
    I-IV-V-I cadence, the trainable relative-pitch skill, rather than
    absolute-key naming, which would require perfect pitch.
\end{itemize}

\subsection{Optional AI features and the local-first exception}
\label{sec:coaching}

The system exposes five optional, opt-in features backed by an external
large-language model, and they are the only places where any data leaves the
machine. They share one design contract, so the exception to local-only
operation stays small and uniform. Each fires only on an explicit button press;
each sends only a compact text or numeric summary, never audio, never a recording,
and not even a file name; and each renders the model's reply as plain text rather
than as markup, so untrusted output cannot inject content into the page. The API
key is read server-side, per request, from a local, git-ignored configuration
file and is never exposed to the browser. All five share a single transport
module, and all are off by default and inert until a key is configured, so the
rest of the system remains strictly local.

The five features are:

\begin{itemize}
    \item \textbf{Sing-Along coaching.} After a take, the client first computes a
    deterministic, human-readable analysis of the performance in the browser: the
    signed pitch bias (consistently sharp or flat), drift from the start of the
    take to the end, octave slips, accuracy of leaps versus stepwise motion, the
    weakest register, and the worst individual notes with their bar numbers. This
    analysis is always shown and needs no network access. Only on request is its
    numeric summary, with the musical context, sent for spoken-language feedback
    and practice suggestions.
    \item \textbf{Practice plan.} On the hub, the aggregate progress report,
    computed locally from the training history, can be sent for a one-week plan
    whose items name the application's own drills. The last plan is cached in the
    browser and keyed by a hash of the report, so opening the hub never issues a
    request on its own.
    \item \textbf{Natural-language editing.} In Manual mode, the ``Ask'' box sends
    a text listing of the melody and the user's instruction, and the model returns
    a small structured edit script; the browser validates and applies it. Common
    instructions are handled entirely offline by a grammar and never reach this
    path.
    \item \textbf{Whole-melody reorganisation.} Also in Manual mode, a
    ``Reorganise'' button sends the current melody's pitches and beat lengths and
    receives a cleaned-up rewrite, with durations snapped to simple note values
    and tied slivers merged; the browser validates it and applies it as one
    undoable step. Where the Ask box makes a bounded local edit, this rewrites the
    whole line for readability.
    \item \textbf{Reharmonisation.} On the result card, a requested style and the
    per-measure pitch-class profile of the melody are sent, and the model proposes
    a new chord progression, which the system validates, spells, scores, and
    engraves locally.
\end{itemize}

The reply requests are handled server-side by synchronous route handlers, which
the web framework runs on a worker thread, so a call of up to a minute does not
block the event loop. The default model is a reasoning model that spends hidden
tokens before its answer, so the structured features (editing, reorganisation,
and reharmonisation) request a larger token budget than the prose features and can
be pointed at a faster model through configuration. Beyond these five calls, HumJob's
operation is entirely local.

\subsection{Tempo detection and its known fragility}
\label{sec:tempo-design}

Because a wrong tempo is the usual cause of tied-note artifacts, the
application offers a ``find my tempo'' button. Its purpose is narrow: the user
hums a few even beats, the system estimates a comfortable tempo, and the user
then records the real take to a click at that tempo, so that quantisation lands
on whole beats.

The estimator computes an onset-strength envelope with
librosa~\cite{librosa}, then finds the dominant periodicity of that envelope by
autocorrelation, weighted by a prior centred on a comfortable humming tempo,
and folds the result into a plausible range by octaves. The important property,
and a known weakness examined in Section~\ref{sec:bpm}, is that this returns a
\emph{single global tempo} averaged over the whole clip, tilted toward the
prior. It assumes a roughly constant tempo and reasonably clean onsets, and it
is fragile when either assumption is violated.

A second, independent safeguard sits in the quantiser itself. Because even a
careful hum drifts a few percent off the metronome, the quantiser refines the
tempo before snapping: it searches a bounded band around the stated tempo (about
$\pm 25\%$) and adopts the beat length that makes the snapped rhythm metrically
simplest, leaving a take that is already on the beat untouched. The score is
still notated at the tempo the user set; only the grid the onsets snap to follows
the performed tempo. This absorbs a moderate tempo error, whether from the user
or from the detector above, but not a gross one, which is why the detector's
fragility still matters for large errors (Section~\ref{sec:bpm}).

\subsection{Technology stack and running the system}
\label{sec:stack}

Table~\ref{tab:stack} lists the main components. The whole system runs locally.
The web application is a FastAPI server that serves a static browser front-end
and reloads on source edits; the command-line tool transcribes a file directly.

\begin{table}[H]
\centering
\caption{Technology stack.}
\label{tab:stack}
\begin{tabular}{l p{0.60\textwidth}}
\toprule
\textbf{Layer} & \textbf{Technology} \\
\midrule
Core pipeline     & Python 3.12, NumPy \\
Pitch backends    & pYIN via librosa~\cite{librosa}; torchcrepe~\cite{crepe};
pesto-pitch~\cite{pesto}; penn / FCNF0++~\cite{penn}; basic-pitch on ONNX
Runtime~\cite{basicpitch} \\
Onset / tempo / chroma & librosa~\cite{librosa} \\
Score model, MIDI, MusicXML & music21~\cite{music21} \\
Engraving         & Verovio~\cite{verovio} (server-side, and a vendored
WebAssembly build in the browser) \\
Web backend       & FastAPI, Uvicorn (with automatic reload) \\
Web front-end     & Static HTML / CSS / JavaScript, Web Audio API \\
AI features (optional) & Shared httpx client to an external, OpenAI-compatible
language-model API (DeepSeek); powers coaching, the practice plan,
natural-language editing, and reharmonisation; off by default, opt-in, text or
numeric summaries only \\
Evaluation        & mir\_eval~\cite{mireval} for note metrics \\
\bottomrule
\end{tabular}
\end{table}

% =========================================================================
% SECTION 4: EVALUATION AND RESULTS
% =========================================================================
\section{Evaluation and Results}

\subsection{Evaluation methodology}

This section evaluates the transcriber, the toolkit's technical core and the only
component that estimates a hidden ground truth. The other five tools are
functional deliverables rather than objects of measurement: their pure logic (the
client score-building and enharmonic spelling, the transposer's key arithmetic,
the sing-along scoring, and the ear-trainer question generation) is covered by
unit tests, and Appendix~A demonstrates each running. There is no accuracy
benchmark for them because they are utilities, not estimators, so the
measurements below concern the transcription pipeline alone.

The project's standard evaluation of the transcriber is a synthetic fixture
harness. It renders a
set of known melodies as audio, runs the pipeline once per fixture, and
compares the output against the known reference. The fixtures cover eight
melodies: a C major scale, an A minor scale, an arpeggio, repeated notes, a
melody with silence, octave leaps, a Twinkle excerpt, and a mixed-rhythm
phrase.

The harness scores two different things, because they measure different
failures:

\begin{itemize}
    \item \textbf{Note detection (pitch and onset F1)}, computed with
    mir\_eval~\cite{mireval}. A predicted note matches a reference note when its
    pitch is correct and its onset is within a tolerance; durations are ignored
    and no offset tolerance is applied. This is run in two flavours: a
    \emph{clean} take that the regression gate expects to pass at F1 $=1.0$, and
    a \emph{realistic} take with wide vibrato, tremolo, pitch drift, jittered
    consonant-gap lengths, and partial consonant closures.
    \item \textbf{Rhythm}, added later, because note F1 is rhythm-blind. This
    scores the quantised onset and duration of each note against an intended
    grid and is described in Section~\ref{sec:results-rhythm}.
\end{itemize}

Both of these are aggregate scores, and both saturated at ceiling during the
project (Sections~\ref{sec:results-note} and~\ref{sec:results-rhythm}). Because a
saturated number can no longer drive improvement and says nothing about real
hums, the evaluation was then widened into four kinds of evidence, presented in
Section~\ref{sec:four-evidence}, each designed to need no precise performer:
diagnostic error-rates on a diversified corpus, metamorphic correctness
properties, a reconstruction round-trip on real hums, and a calibration of the
synthetic realism against real takes. The saturated F1 results below are the
baseline that motivated that widening.

\subsection{The ground-truth problem}
\label{sec:groundtruth}

Before the results, their central limitation must be stated plainly. \textbf{All
of the ground truth used in this report is synthetic.} The fixtures are
melodies the project itself renders to audio, so their intended pitch, onset,
and duration are known exactly by construction. This is a genuine strength for
regression testing, but it is a genuine weakness for validity: a synthetic hum
is more regular than a real one, and a benchmark built only on synthetic hums
can be saturated while the system still fails on real input.

The project has no labelled corpus of real hummed melodies. The reason is
specific to this task and is the same kind of problem that limits any
transcription benchmark: obtaining ground truth for a real hum requires knowing
the exact pitch, onset, and duration of every note the person intended, and a
person cannot reliably label the pitch of their own hum by ear. As a result,
every note-detection and rhythm figure in this report measures agreement with
\emph{synthetic} ground truth and must be read as a regression signal, not as a
measure of real-world accuracy. Section~\ref{sec:eval-limits} proposes a
concrete remedy.

\subsection{Note-detection results}
\label{sec:results-note}

On the clean fixtures the pipeline scores a mean note F1 of $1.000$. The more
informative number is the realistic take, where the segmenter historically
struggled. Table~\ref{tab:note-progression} records the mean realistic F1
across three engineering iterations. The first raised it from $0.799$ to
$0.931$ by making segmentation and consolidation grid-aware; the second raised
it to $1.000$ by adding the spectral octave correction of
Section~\ref{sec:octave}, which fixed the last failing fixture (the octave
leaps, previously returned an octave flat).

\begin{table}[H]
\centering
\caption{Mean note F1 on the realistic fixtures across three iterations. The
realistic take is the expressive one (wide vibrato, tremolo, drift, jittered
consonant gaps, partial closures).}
\label{tab:note-progression}
\begin{tabular}{l l l}
\toprule
\textbf{Iteration} & \textbf{Change} & \textbf{Mean realistic F1} \\
\midrule
Baseline              & coarse valley splitter                 & $0.799$ \\
Grid-aware            & grid-aware segment + consolidate        & $0.931$ \\
Octave correction     & spectral subharmonic repair             & $1.000$ \\
\bottomrule
\end{tabular}
\end{table}

Table~\ref{tab:note-perfixture} gives the per-fixture note metrics on the
realistic take at the current state. Every fixture scores $1.000$ on precision,
recall, and F1, and the key is detected correctly in every case. This is the
saturation referred to throughout this report: the note-detection half of the
synthetic harness no longer discriminates between a better and a worse
segmenter, so it can no longer drive improvement on its own.

\begin{table}[H]
\centering
\caption{Per-fixture note metrics on the realistic take, current state. P,
R, F1 are precision, recall, and F1 from mir\_eval; ref/est is
reference versus detected note counts.}
\label{tab:note-perfixture}
\begin{tabular}{l l l l l l}
\toprule
\textbf{Fixture} & \textbf{P} & \textbf{R} & \textbf{F1} & \textbf{Key} &
\textbf{ref/est} \\
\midrule
c\_major\_scale  & $1.00$ & $1.00$ & $1.00$ & C major & 8/8 \\
a\_minor\_scale  & $1.00$ & $1.00$ & $1.00$ & A minor & 8/8 \\
arpeggio         & $1.00$ & $1.00$ & $1.00$ & C major & 7/7 \\
repeated\_notes  & $1.00$ & $1.00$ & $1.00$ & C major & 5/5 \\
with\_silence    & $1.00$ & $1.00$ & $1.00$ & C major & 5/5 \\
octave\_leaps    & $1.00$ & $1.00$ & $1.00$ & C major & 5/5 \\
twinkle          & $1.00$ & $1.00$ & $1.00$ & C major & 14/14 \\
mixed\_rhythm    & $1.00$ & $1.00$ & $1.00$ & C major & 5/5 \\
\midrule
\textbf{mean}    &        &        & $1.000$ &        & \\
\bottomrule
\end{tabular}
\end{table}

\subsection{The rhythm evaluation}
\label{sec:results-rhythm}

The note F1 above is rhythm-blind by construction: it ignores durations and
accepts any onset within a tolerance, so it says nothing about whether the
written rhythm is correct. This matters because the two failures reported most
often on real hums are wrong rhythm (durations and onsets on the wrong beats)
and over-splitting (one note returned as tied slivers), and neither is visible
to note F1. A separate rhythm metric was therefore added. It compares the
quantised onset and duration of each note against an \emph{intended grid} read
straight from each fixture's specification, and reports the fraction of notes
with the right onset, the right duration, and both. The headline number is
``both'': right beat and right written length.

Two conditions were run, chosen to match the two things that actually break
real hums. Table~\ref{tab:rhythm} gives the per-fixture results.

\begin{table}[H]
\centering
\caption{Rhythm accuracy under two conditions. ``on'', ``dur'', and ``both''
are the fraction of notes with the correct onset, correct duration, and both;
``err'' is the mean onset error in milliseconds. The jitter condition applies
human timing jitter ($\pm 30$ ms) at the correct tempo; the wrong-tempo
condition applies a $+5\%$ tempo error to on-grid timing.}
\label{tab:rhythm}
\begin{tabular}{l r r r r @{\hskip 1.5em} r r r r}
\toprule
 & \multicolumn{4}{c}{\textbf{Timing jitter ($\pm 30$ ms)}}
 & \multicolumn{4}{c}{\textbf{Wrong tempo ($+5\%$)}} \\
\cmidrule(lr){2-5}\cmidrule(lr){6-9}
\textbf{Fixture} & \textbf{on} & \textbf{dur} & \textbf{both} & \textbf{err} &
\textbf{on} & \textbf{dur} & \textbf{both} & \textbf{err} \\
\midrule
c\_major\_scale  & $1.00$ & $1.00$ & $1.00$ & $0$ & $1.00$ & $1.00$ & $1.00$ & $0$ \\
a\_minor\_scale  & $1.00$ & $1.00$ & $1.00$ & $0$ & $1.00$ & $1.00$ & $1.00$ & $0$ \\
arpeggio         & $1.00$ & $1.00$ & $1.00$ & $0$ & $1.00$ & $1.00$ & $1.00$ & $0$ \\
repeated\_notes  & $1.00$ & $1.00$ & $1.00$ & $0$ & $1.00$ & $1.00$ & $1.00$ & $0$ \\
with\_silence    & $1.00$ & $0.80$ & $0.80$ & $0$ & $1.00$ & $1.00$ & $1.00$ & $0$ \\
octave\_leaps    & $1.00$ & $1.00$ & $1.00$ & $0$ & $1.00$ & $1.00$ & $1.00$ & $0$ \\
twinkle          & $1.00$ & $1.00$ & $1.00$ & $0$ & $1.00$ & $1.00$ & $1.00$ & $0$ \\
mixed\_rhythm    & $1.00$ & $1.00$ & $1.00$ & $0$ & $1.00$ & $1.00$ & $1.00$ & $0$ \\
\midrule
\textbf{mean}    &        &        & $\mathbf{0.975}$ & & & & $\mathbf{1.000}$ & \\
\bottomrule
\end{tabular}
\end{table}

Two findings stand out. First, on-grid at the correct tempo the quantiser is
perfect ($1.000$), and under human timing jitter the mean stays at $0.975$: every
fixture but one scores perfectly, and only the silence fixture loses a written
duration, so the quantiser is robust to the timing drift of a real performance.
Second, a $5\%$ tempo error is also recovered to $1.000$, because the quantiser
refines the tempo before snapping (Section~\ref{sec:tempo-design}): it searches a
bounded band around the stated tempo and adopts the beat length that makes the
snapped rhythm metrically simplest, so a hum that drifts a few percent off the
metronome is pulled back onto whole beats instead of into the off-beat onsets
that notation renders as tied slivers. Within that band the rhythm half of the
synthetic harness is therefore saturated as well.

Figure~\ref{fig:rhythm} plots the two conditions side by side.

\begin{figure}[H]
\centering
\reportfig{0.85}{eval-rhythm.png}{Rhythm ``both'' accuracy per fixture under
human timing jitter (correct tempo) and under a $5\%$ tempo error. Both
conditions are at or near ceiling: the tempo-refinement stage absorbs the tempo
error, and only the silence fixture drops slightly under jitter.}
\caption{Rhythm accuracy under timing jitter and under a wrong tempo.}
\label{fig:rhythm}
\end{figure}

\subsection{Beyond one saturated number: four kinds of evidence}
\label{sec:four-evidence}

A single F1 pinned at ceiling cannot separate a better segmenter from a worse
one, and it is entirely synthetic, so it proves nothing about real hums. The
usual fix, a labelled corpus of real hums, is blocked here for the reason in
Section~\ref{sec:groundtruth}: the author is not a precise singer and no labelled
real-hum dataset exists. The evaluation was therefore widened into four kinds of
evidence, each chosen so that it needs no precise performer, and each reported
below with an explicit statement of what it does and does not prove. Together
they replace one saturated number with a picture that stays useful.

\subsubsection{Evidence 1: diagnostic error-rates on a diversified corpus}
\label{sec:diag}

The first step is to stop hiding failures inside one aggregate. Instead of note
F1, the pipeline is scored on a set of diagnostic error-rates that each pick out
one failure mode, over a larger, more varied corpus. The corpus has sixteen
melodies: the original eight fixtures, three public-domain tunes (Ode to Joy,
Mary Had a Little Lamb, Frere Jacques), and five procedurally generated melodies
(in-key random walks, an arpeggio, an interval study, and a varied-rhythm line),
all rendered with known ground truth. Each melody is rendered at three realism
profiles: \emph{clean}, \emph{realistic}, and a new, harder \emph{hard} profile
that adds amplitude shimmer, pitch creak, a breathy noise floor, a short room
reverberation, and pink microphone noise on top of the realistic vibrato and
partial closures.

The error-rates are: the signed over/under-split rate (estimated minus reference
note count, over the reference count, so a positive value is over-splitting and a
negative value is merging or dropping notes); the octave-error rate (aligned
notes off by an exact octave); the same-pitch repeat recall (of the reference's
adjacent same-pitch pairs, the fraction that survive as two distinct notes, which
is the ``merged repeats'' complaint); and the off-grid rate (notes whose
quantised onset or duration leaves the simple-note grid, which is the
``tied-sliver'' complaint). Table~\ref{tab:diag} gives the corpus means per
profile.

\begin{table}[H]
\centering
\caption{Diagnostic error-rates, corpus means over the sixteen melodies, at three
realism profiles. F1 is the mean note F1 for reference. ``split'' is the signed
over/under-split rate (negative is merging or dropping). ``octave'' is the
octave-error rate. ``repeat'' is the same-pitch repeat recall (higher is better).
``off-grid'' is the tied-sliver / off-grid rate. No melody collapsed to zero
notes at any profile.}
\label{tab:diag}
\begin{tabular}{l r r r r r}
\toprule
\textbf{Profile} & \textbf{F1} & \textbf{split} & \textbf{octave} &
\textbf{repeat} & \textbf{off-grid} \\
\midrule
Clean      & $1.000$ & $+0.00$ & $0.00$ & $1.00$ & $0.00$ \\
Realistic  & $0.979$ & $-0.04$ & $0.00$ & $0.61$ & $0.00$ \\
Hard       & $0.918$ & $-0.04$ & $0.00$ & $0.63$ & $0.00$ \\
\bottomrule
\end{tabular}
\end{table}

Three things stand out. First, the corpus discriminates again: the numbers are no
longer all $1.000$, and they order the three profiles as expected (clean best,
hard worst on F1). Second, the most useful single number is the repeat recall,
which falls from $1.00$ on clean input to about $0.61$ on realistic input: on an
expressive take the segmenter merges roughly four in ten of the same-pitch
repeats it should keep apart. That is the long-suspected ``merged repeats''
weakness, now measured rather than only reported informally. Third, the octave
rate stays at zero (the spectral octave correction holds) and the off-grid rate
stays at zero at the correct tempo (the quantiser keeps notes on simple values).

One expectation was not met, and the report states it plainly. The \emph{hard}
profile was designed to provoke \emph{over-splitting} (a held note shattered into
slivers), the failure the too-regular synthesiser could not previously show.
Direct measurement showed the opposite. This pipeline does not over-split even
under heavy perturbation; its median-smoothing and consolidation stages hold, and
hard input instead degrades by \emph{merging or dropping} notes. Pushing the
synthesiser far enough to fragment a note instead silences it. The signed split
rate is therefore mildly negative, never positive. This is a genuine finding
about the pipeline's behaviour, and it corrects the earlier expectation recorded
in this report.

\textbf{What this proves:} the pipeline's behaviour on varied, expressive input is
now measured across distinct failure modes, and the corpus can once again tell a
better segmenter from a worse one. The dominant expressive-input failure is
localised to merged same-pitch repeats, not to octave errors or off-grid rhythm.
\textbf{What it does not prove:} the ground truth is still synthetic, so the
absolute rates depend on how faithful the realism profiles are, which is the
subject of Evidence~4.

\subsubsection{Evidence 2: metamorphic correctness properties}
\label{sec:meta}

The second kind of evidence needs no reference melody at all. A metamorphic test
asserts a \emph{relation} that must hold between a take and a transformed version
of the same take, so the ``label'' is the transformation, not a known pitch. Five
properties are checked: transpose invariance (shifting the pitch by $k$ semitones
shifts every output note by $k$ and leaves the count unchanged); gain invariance
(scaling the amplitude leaves the transcription identical); additive-noise
robustness (adding white noise down to a $30$~dB signal-to-noise ratio leaves the
notes unchanged); silence-pad invariance (prepending or appending silence leaves
the same notes at the same values); and tempo, or time-scale, invariance
(performing the same melody at a scaled tempo and reading it at that tempo leaves
the note values unchanged). All five hold on the synthetic fixtures. Transpose
invariance and gain invariance also hold on the real recorded takes, where they
are true regardless of how accurate the hum was.

\textbf{What this proves:} a body of correctness relations the transcriber must
obey, verifiable with no precise performer and partly on real audio; a change that
broke transpose or tempo handling would be caught here. \textbf{What it does not
prove:} an invariance can hold while the transcription is wrong in a way the
transformation also preserves, so these bound whole classes of error but do not
measure absolute accuracy.

\subsubsection{Evidence 3: a reconstruction round-trip on real hums}
\label{sec:roundtrip}

The third kind of evidence is the only one computed on real voice. Because a real
hum has no ground truth, the transcription is scored against \emph{its own audio}:
the transcribed notes are resynthesised to a plain tone track, and that
reconstruction is compared with the original recording. Pitch agreement is the
fraction of jointly-voiced frames within $50$ cents; onset agreement is an
onset-time F1 between the audio's detected onsets and the transcribed onsets; and
a single ``explains-audio'' summary averages the two. Table~\ref{tab:roundtrip}
gives the result on the four real takes recorded.

\begin{table}[H]
\centering
\caption{Reconstruction round-trip on the four real recorded takes. ``explains''
is the summary score, ``pitch'' is the fraction of jointly-voiced frames within
$50$ cents, ``onset'' is the onset-time F1, and ``cents'' is the median absolute
pitch error over jointly-voiced frames. Higher is better except for cents.}
\label{tab:roundtrip}
\begin{tabular}{l r r r r}
\toprule
\textbf{Take} & \textbf{explains} & \textbf{pitch} & \textbf{onset} &
\textbf{cents} \\
\midrule
mary\_1     & $0.59$ & $0.70$ & $0.49$ & $40$ \\
repeated\_1 & $0.75$ & $0.92$ & $0.57$ & $20$ \\
scale\_1    & $0.54$ & $0.73$ & $0.35$ & $20$ \\
twinkle\_1  & $0.61$ & $0.84$ & $0.39$ & $20$ \\
\midrule
\textbf{mean} & $\mathbf{0.62}$ & & & \\
\bottomrule
\end{tabular}
\end{table}

Pitch agreement is high (between $0.70$ and $0.92$), which says the resynthesised
melody tracks the pitch the person actually hummed and that no gross octave slip
occurred. Onset agreement is lower (between $0.35$ and $0.57$), largely because
the onset detector fires on more places in the raw hum than there are notes, so
the number understates true agreement. Crucially, because the score compares each
transcription with \emph{its own} audio and never with an intended melody, an
imprecise hum cannot count against the system: a wrong note the pipeline captures
faithfully still agrees. That is what lets these numbers come from a hummer who is
not a trained singer. The takes are few (four), so they are read as a sanity
signal, not a trend.

One take is worth a separate note. On the Mary-Had-a-Little-Lamb take the pipeline
\emph{over-split} the melody (it returned roughly twice as many notes as intended).
That is the failure the synthetic \emph{hard} profile could not reproduce
(Section~\ref{sec:diag}), so real voice evidently carries cues, or capture noise,
that the synthesiser does not. That take was also quiet and of uncertain tempo, so
it is a single weak data point rather than a measured rate; but it is a concrete
reason not to over-read the synthetic-only finding that this pipeline does not
over-split, and another argument for real ground truth.

\textbf{What this proves:} an unsupervised, real-voice sanity signal that needs no
label and that catches gross pitch, onset, and segmentation mismatch, including
octave slips. \textbf{What it does not prove:} it is a sanity signal, not an
accuracy number. A self-consistent error, such as a constant tuning offset within
$50$ cents, or any error the resynthesis reproduces, leaves the two contours
agreeing and is not caught.

\subsubsection{Evidence 4: calibrating the synthetic realism from real takes}
\label{sec:calib}

The fourth kind of evidence closes the loop back to the first. The diagnostic
error-rates of Evidence~1 are only as honest as the realism profiles they are
measured at, so those profiles are checked against real voice. For each real take
the tracker is run and the quantities the realism knobs model are measured:
vibrato depth and rate, pitch drift, timing jitter against the metronome, the
depth of the consonant energy dips, and amplitude shimmer.
Table~\ref{tab:calib} gives the medians over the four real takes, next to the
value the \emph{realistic} profile assumes.

\begin{table}[H]
\centering
\caption{Realism quantities measured from the four real takes (medians), next to
the value the realistic synthetic profile assumes. The timing-jitter row is the
exception: the realistic pitch profile itself keeps timing clean, so the $30$~ms
shown is the jitter the rhythm test injects (Section~\ref{sec:results-rhythm}),
which is what the measured value is compared against.}
\label{tab:calib}
\begin{tabular}{l r r}
\toprule
\textbf{Quantity} & \textbf{Measured (median)} & \textbf{Realistic profile} \\
\midrule
Vibrato depth (cents) & $26$  & $55$ \\
Vibrato rate (Hz)     & $5.2$ & $5.3$ \\
Pitch drift (cents)   & $18$  & $18$ \\
Timing jitter (ms)    & $131$ & $30$ \\
Closure depth (dB)    & $5$   & $16$ \\
Shimmer (dB)          & $1.1$ & --- \\
\bottomrule
\end{tabular}
\end{table}

The vibrato rate and the pitch drift the synthesiser assumes are close to what
real humming shows: the rate matches almost exactly ($5.2$ against $5.3$~Hz), and
the drift lands on the assumed value ($18$ cents in both). The assumed vibrato
depth is wider than measured ($55$ against $26$ cents), so the realistic profile
is, if anything, harder than a real hum on pitch. The measured timing jitter is
much larger than the value used in the rhythm test ($131$ against $30$~ms), which
says real humming is far looser in time than the rhythm test assumes and is a
reason the tempo work of Section~\ref{sec:bpm} matters.

\textbf{What this proves:} the synthetic realism is grounded in measured numbers
rather than guessed, and two of the central knobs (vibrato rate and drift) match
real voice well. \textbf{What it does not prove:} four takes is a small sample,
and the measurements are estimates from the tracker, not exact.

\subsection{Interpreting the scores}
\label{sec:interpret}

The evidence must be read together and read carefully. A saturated note F1 of
$1.000$ does not mean the system transcribes real hums correctly; it means the
synthetic note-detection benchmark can no longer tell a good segmenter from a bad
one. The four kinds of evidence above were built precisely because of that. The
diagnostic error-rates (Evidence~1) restore a benchmark that discriminates and
localise the dominant expressive-input failure to merged same-pitch repeats; the
metamorphic properties (Evidence~2) give correctness relations that need no
performer; the reconstruction round-trip (Evidence~3) gives a first, unsupervised
real-voice signal; and the realism calibration (Evidence~4) keeps the synthetic
profiles honest. None of these is a real-world accuracy number, because that still
requires real ground truth, but together they are a far more useful picture than
one number at ceiling.

The rhythm evaluation once localised the dominant real-world failure to a single
cause: sensitivity to tempo error. That diagnosis drove the tempo-refinement
stage, which now absorbs a moderate tempo error and, with it, most of the ``wrong
rhythm'' behaviour that a small tempo mismatch used to cause. What remains open is
the tempo detector's behaviour on real humming that drifts beyond the refinement
band (Section~\ref{sec:bpm}), and, more fundamentally, the absence of any real
ground truth to measure absolute accuracy against.

\subsection{The tempo-detection weakness}
\label{sec:bpm}

The rhythm result above shows that a moderate tempo error is now absorbed by the
refinement stage in the quantiser. That stage spans only a bounded band (about
$\pm 25\%$ of the stated tempo), however, so a gross error from the system's own
tempo detector still escapes it, and it is worth stating why that detector is
fragile.

As described in Section~\ref{sec:tempo-design}, the detector returns a single
global tempo: it takes the autocorrelation of the onset-strength envelope over
the whole clip, weights it by a prior centred near a comfortable humming tempo,
and returns the dominant period. This has three consequences on real input.
First, it assumes a constant tempo; a hum that speeds up or slows down has onset
energy at more than one period, and the autocorrelation returns a blurry
compromise rather than either true tempo. Second, the prior actively biases the
estimate toward the centre of its range, so a genuinely slow hum is pulled
upward. Third, it relies on clean, evenly spaced onsets, which real legato or
noisy humming does not always provide. In informal testing, a hum of a roughly
$76$ beat-per-minute melody that included a faster passage was reported as about
$101$, a value suspiciously close to the prior, which is the signature of the
prior dominating a weak autocorrelation. An error of that size (about $+33\%$)
falls outside the refinement band and so is not recovered, which is exactly the
case where the detector's fragility still bites. This weakness is the leading
item of future work in Section~\ref{sec:directions}.

% =========================================================================
% SECTION 5: LIMITATIONS AND FUTURE WORK
% =========================================================================
\section{Limitations and Future Work}

\subsection{Evaluation limitations}
\label{sec:eval-limits}

The main limitation is the one stated in Section~\ref{sec:groundtruth}: no
labelled corpus of real hummed melodies exists, so no figure in this report is a
real-world accuracy number. The aggregate note and rhythm F1 both saturated
during the project. The response was to widen the evaluation
(Section~\ref{sec:four-evidence}) rather than to keep reporting one number at
ceiling: the diagnostic error-rates over the diversified corpus discriminate
again, the metamorphic properties and the reconstruction round-trip add signals
that need no precise performer, and the realism calibration keeps the synthetic
profiles honest. What none of these gives is absolute accuracy on real voice,
which still needs real ground truth. A second finding narrows the earlier plan:
the over-splitting failure was expected to appear once shimmer, breath, and creak
were added to the synthesiser, but direct measurement showed this pipeline does
not over-split; its smoothing and consolidation defences hold, and hard input
degrades by merging or dropping notes instead. So the open evaluation gap is not
over-split synthesis but real ground truth.

The recommended remedy is not a larger synthetic set but a small set of real,
labelled hums. The obstacle is ground truth, since a person cannot label the
pitch of their own hum by ear. The plan, refined during this project, inverts
the labelling. Rather than hum first and label afterward, the label is fixed
before the recording: the user hums a melody they know, to a click at a known
tempo, and the rhythm ground truth comes from the known tune and the click
grid. The pitch of each note is confirmed with an external pitch-finder
application used as an assistive aid, with the user verifying each note against
the tune they knowingly hummed, rather than trusting the external tool as an
oracle. This is important because an external pitch finder is itself another
pitch detector; if it is trusted blindly and it happens to share an error with
the pipeline, the evaluation would score that shared error as correct. Keeping
the user in the loop, and drawing the rhythm truth from the click rather than
from the tool, avoids that circularity. The resulting verified recordings, with
their known notes and tempo, would give the project its first real ground truth.

\subsection{System deficiencies}

Several deficiencies of the current system are known and tracked.
Table~\ref{tab:deficiencies} lists them together with the planned remedy for
each.

\begin{longtable}{>{\raggedright\arraybackslash}p{0.30\textwidth} >{\raggedright\arraybackslash}p{0.58\textwidth}}
\caption{Known deficiencies of the current system and planned work.}
\label{tab:deficiencies} \\
\toprule
\textbf{Area} & \textbf{Description and planned remedy} \\
\midrule
\endfirsthead
\toprule
\textbf{Area} & \textbf{Description and planned remedy} \\
\midrule
\endhead
\midrule
\multicolumn{2}{r}{\footnotesize\itshape Continued on the next page.} \\
\endfoot
\bottomrule
\endlastfoot
Tempo detection fragility & The ``find my tempo'' estimator returns a single
global tempo biased by a prior, and is unreliable on real, variable-tempo,
slow, or legato humming (Section~\ref{sec:bpm}); the quantiser's tempo-refinement
stage absorbs a moderate error but not a gross one. Planned remedy: flatten or
remove the prior, and build the onset envelope from the pipeline's own voicing
onsets rather than raw spectral flux. \\
\addlinespace
Merged same-pitch repeats & On expressive input the segmenter merges roughly four
in ten of the same-pitch repeats it should keep apart (repeat recall about $0.61$,
Section~\ref{sec:diag}). This, not over-splitting, is the dominant expressive-input
failure. Planned remedy: strengthen the same-pitch onset cue and re-measure repeat
recall as the target. \\
\addlinespace
Over-split does not occur, but was expected & The over-split failure (one note
returned as slivers) was expected to appear once shimmer, breath, and creak were
added to the synthesiser. Measurement showed this pipeline does not over-split; it
merges or drops notes instead (Section~\ref{sec:diag}). No remedy needed; the
expectation is corrected. \\
\addlinespace
No real ground truth & No figure is a real-world accuracy number
(Section~\ref{sec:groundtruth}); the four real takes feed only the round-trip and
calibration, which are performer-tolerant, not an accuracy score. Planned remedy:
the hum-to-a-known-score capture flow of Section~\ref{sec:eval-limits}. \\
\addlinespace
Aggregate F1 saturated (addressed) & The note-detection and rhythm F1 are at
ceiling on the fixtures and no longer discriminate. Addressed by the diagnostic
error-rates, metamorphic properties, round-trip, and realism calibration of
Section~\ref{sec:four-evidence}; the remaining gap is real ground truth. \\
\addlinespace
Monophonic only & The transcription pipeline assumes a single voice, one note
at a time. Polyphonic or chordal input is out of scope; the Pitch Finder tab
uses a separate chroma-based method for full songs. \\
\addlinespace
Optional backend availability & The most precise backends (PESTO, FCNF0++) are
optional installs with heavier dependency trees, and FCNF0++ downloads weights
on first use. The pipeline falls back to pyin when they are absent, at a
precision cost. \\
\addlinespace
Fixed metre assumption & Quantisation assumes a simple metre and a known,
steady tempo. Time-signature changes and expressive rubato are not modelled. \\
\end{longtable}

\subsection{Future directions}
\label{sec:directions}

Three directions are recorded for the work after this report, in priority
order.

\begin{itemize}
    \item \textbf{(A) Tempo robustness (leading candidate).} With the quantiser's
    tempo-refinement stage now absorbing moderate errors, the tempo detector is
    the largest remaining lever on real-world rhythm accuracy for the gross errors
    that fall outside the refinement band (Section~\ref{sec:bpm}), and it is
    directly testable because the fixtures have known tempos. Two cheap,
    high-impact changes: flatten or remove the tempo prior so slow tempos are not
    pulled upward, and build the onset envelope from the pipeline's own consonant
    onsets rather than from raw spectral flux.
    \item \textbf{(B) Real-hum ground truth.} The ultimate unlock, blocked only
    on the labelling method, which Section~\ref{sec:eval-limits} resolves. Four
    real takes already feed the reconstruction round-trip and the realism
    calibration (Sections~\ref{sec:roundtrip} and~\ref{sec:calib}); a larger
    labelled set, captured by the same method, would turn those from a sanity
    signal into a firmer one and let the project measure real-world accuracy for
    the first time.
    \item \textbf{(C) The merged-repeat cue.} With over-splitting shown not to
    occur, the measured weakness is the opposite one: merged same-pitch repeats on
    expressive input (Section~\ref{sec:diag}). Strengthening the same-pitch onset
    cue, with repeat recall as the target metric, is the concrete next step on
    segmentation.
\end{itemize}

\subsection{Towards production readiness}

HumJob is a research and academic project, not a production service. It runs
locally, which is a deliberate design choice, but a wider deployment would
additionally require packaging the optional neural backends so they install
reliably, a broader automated test suite over real recordings, and error
handling for the many ways a real microphone recording can be unusable (clipping,
background noise, too quiet). None of these is conceptually hard, but together
they are a meaningful body of work beyond the scope of this course project.

% =========================================================================
% SECTION 6: CONCLUSION
% =========================================================================
\section{Conclusion}

This project set out to transcribe a hummed melody into notes, key, and chords,
by accepting two small constraints on how the user hums and building the whole
pipeline around them. The transcriber that resulted is a local-first pipeline of
pure functions over three data types, with five interchangeable pitch backends, a
grid-aware segmenter and consolidator, and a spectral octave-correction stage.
Around that core the project grew into a local-first toolkit for the singing
voice, delivered as a web application of seven tabs: the transcriber and its
interactive editor, a pitch finder, a transposer, a real-time monitor and vocal
trainer, a sing-along trainer, an ear trainer, and a practice hub. All of it runs
on the user's own machine; the only deliberate exception is a small family of
five opt-in, text-only AI features (coaching, a practice plan, natural-language
editing, whole-melody reorganisation, and reharmonisation), each of which sends a
compact summary and never any audio.

With respect to the research questions, the honest verdicts are mixed.
\textbf{RQ1} is supported by construction and by the evaluation: the two
interface constraints do turn the hard parts into manageable ones, and the
grid-aware stages that exploit them are what carried the realistic note F1 from
$0.799$ to $1.000$. \textbf{RQ2} is supported: the diagnosed failures were
segmentation and quantisation failures over an essentially correct pitch
contour, not pitch-model failures. The clearest case is the octave error, which
was proven to be a Viterbi decoding artifact rather than a fault in the audio,
and was repaired by a spectral test rather than by changing the model.
\textbf{RQ3} produced the most important lesson of the project. The aggregate
note and rhythm F1 the system was tuned against both saturated at $1.000$ while
the system is still not reliable on real hums, because a synthetic hum is more
regular than a real one and no labelled real corpus exists to measure against.
Rather than keep reporting a number at ceiling, the evaluation was rebuilt around
four kinds of evidence that need no precise singer
(Section~\ref{sec:four-evidence}). These discriminate again and, between them,
say something concrete: the dominant expressive-input failure is merged
same-pitch repeats (not, as expected, over-splitting, which this pipeline does
not do); the transcriber obeys a set of correctness invariances that also hold on
real audio; and a real hum's transcription can be checked against its own audio
as an unsupervised sanity signal. What none of them gives is absolute real-world
accuracy, which still requires real ground truth.

The main deliverable is therefore not a headline accuracy number. It is a
working, well-structured toolkit whose transcription core is paired with an
honest evaluation that says clearly where that core stands: strong on the
aggregate synthetic benchmark, discriminated again by richer diagnostics, and not
yet validated on real input. Four real takes already anchor the round-trip and
realism-calibration signals; the most valuable next step is to grow that into a
larger, trustworthy set of real, labelled hums, so that the system's real-world
accuracy can be measured directly rather than inferred from a synthetic proxy.

% =========================================================================
% REFERENCES
% =========================================================================
\newpage
\addcontentsline{toc}{section}{References}
\begin{thebibliography}{99}

\bibitem{pyin} M. Mauch and S. Dixon. \emph{pYIN: A fundamental frequency
estimator using probabilistic threshold distributions.} IEEE International
Conference on Acoustics, Speech and Signal Processing (ICASSP), 2014.

\bibitem{crepe} J. W. Kim, J. Salamon, P. Li, and J. P. Bello. \emph{CREPE: A
Convolutional Representation for Pitch Estimation.} ICASSP, 2018.
arXiv:1802.06182. Implementation: torchcrepe,
\url{https://github.com/maxrmorrison/torchcrepe}.

\bibitem{pesto} A. Riou, S. Lattner, G. Hadjeres, and G. Peeters. \emph{PESTO:
Pitch Estimation with Self-supervised Transposition-equivariant Objective.}
International Society for Music Information Retrieval Conference (ISMIR), 2023.
\url{https://github.com/SonyCSLParis/pesto}.

\bibitem{penn} M. Morrison, C. Hsieh, N. Pruyne, and B. Pardo.
\emph{Cross-domain Neural Pitch and Periodicity Estimation.} 2023.
arXiv:2301.12258. Implementation: penn,
\url{https://github.com/interactiveaudiolab/penn}.

\bibitem{basicpitch} R. M. Bittner, J. J. Bosch, D. Rubinstein,
G. Meseguer-Brocal, and S. Ewert. \emph{A Lightweight Instrument-Agnostic Model
for Polyphonic Note Transcription and Multipitch Estimation.} ICASSP, 2022.
\url{https://github.com/spotify/basic-pitch}.

\bibitem{krumhansl} C. L. Krumhansl. \emph{Cognitive Foundations of Musical
Pitch.} Oxford University Press, 1990. (Krumhansl-Kessler / Krumhansl-Schmuckler
probe-tone key profiles.)

\bibitem{librosa} B. McFee, C. Raffel, D. Liang, D. P. W. Ellis, M. McVicar,
E. Battenberg, and O. Nieto. \emph{librosa: Audio and Music Signal Analysis in
Python.} Proceedings of the 14th Python in Science Conference (SciPy), 2015.
\url{https://librosa.org/}.

\bibitem{music21} M. S. Cuthbert and C. Ariza. \emph{music21: A Toolkit for
Computer-Aided Musicology and Symbolic Music Data.} ISMIR, 2010.
\url{https://web.mit.edu/music21/}.

\bibitem{verovio} L. Pugin, R. Zitellini, and P. Roland. \emph{Verovio: A
library for Engraving MEI Music Notation into SVG.} ISMIR, 2014.
\url{https://www.verovio.org/}.

\bibitem{mireval} C. Raffel, B. McFee, E. J. Humphrey, J. Salamon, O. Nieto,
D. Liang, and D. P. W. Ellis. \emph{mir\_eval: A Transparent Implementation of
Common MIR Metrics.} ISMIR, 2014. \url{https://github.com/craffel/mir_eval}.

\bibitem{hollien} Hollien, Dew, and Philips. \emph{Phonational Frequency Ranges
of Adults.} Journal of Speech and Hearing Research, 14:755, 1971. (Population
norm for the vocal range finder: means of 38 semitones for males and 37 for
females; the standard deviation used is a stated assumption, not from this
study.)

\end{thebibliography}

% =========================================================================
% APPENDIX A - SYSTEM DEMONSTRATION (SCREENSHOTS)
% =========================================================================
\newpage
\appendix
\addcontentsline{toc}{section}{Appendix A: System demonstration}
\section*{Appendix A: System demonstration}
\renewcommand{\thesubsection}{A.\arabic{subsection}}
\setcounter{subsection}{0}

This appendix shows the running web application from a user's perspective. Each
subsection corresponds to one tab. The screenshots are captured against a local
instance of the stack; small visual details depend on the user's browser and
operating system.

\subsection{Hub}

The Hub is the home tab: a read-only practice dashboard that folds the history the
other tabs record locally into offline trend lines and a compact progress report.
Its Coach card adds the opt-in ``Get a practice plan'' button.

\begin{figure}[H]
\centering
\reportfig{0.85}{screenshot-hub.png}{The Hub tab: offline progress trends folded
from the local training history, and the opt-in ``Get a practice plan'' card.}
\caption{Hub practice dashboard.}
\label{fig:scr-hub}
\end{figure}

\subsection{Transcriber: Auto mode}

The Transcriber tab records or uploads a hum and returns the automatic draft:
the engraved sheet music, the detected key, and the suggested chords, with
downloads for MIDI and MusicXML. The chords can be reworked in place: clicking a
chord offers offline alternatives, and a ``Reharmonise'' row asks the opt-in
language model for a new progression in a requested style.

\begin{figure}[H]
\centering
\reportfig{0.85}{screenshot-transcriber.png}{The Transcriber tab in Auto mode:
the engraved server-rendered sheet, the detected key, and the export buttons.}
\caption{Transcriber: Auto mode.}
\label{fig:scr-transcriber}
\end{figure}

\subsection{Transcriber: Manual mode}

Manual mode is the in-browser staff editor for correcting the draft. The
reference pitch strip shows the hummed contour against the chosen notes, and the
editing toolbar offers pitch, duration, merge, split, delete, and insert with
undo and redo. An ``Ask'' box also accepts plain-language edits, handled by an
offline grammar for common phrasings and otherwise by the opt-in language model.

\begin{figure}[H]
\centering
\reportfig{0.85}{screenshot-manual.png}{The Transcriber tab in Manual mode: the
editable staff, the reference pitch strip, and the editing toolbar.}
\caption{Transcriber: Manual mode.}
\label{fig:scr-manual}
\end{figure}

\subsection{Pitch Finder}

The Pitch Finder tab analyses an uploaded clip for key, tempo, and a Camelot
code, using the chroma-based method that works on polyphonic audio.

\begin{figure}[H]
\centering
\reportfig{0.85}{screenshot-pitchfinder.png}{The Pitch Finder tab: detected
key, tempo, Camelot code, and technical statistics for an uploaded clip.}
\caption{Pitch Finder.}
\label{fig:scr-pitchfinder}
\end{figure}

\subsection{Realtime and vocal trainer}

The Realtime tab is a client-side pitch monitor and guitar tuner with
vocal-training drills: a target-note lane, a match game, a scale trainer,
vibrato analysis, a vocal range finder, and a per-session progress panel.

\begin{figure}[H]
\centering
\reportfig{0.85}{screenshot-realtime.png}{The Realtime tab: the live voice
monitor with the practice sub-mode selector (Free / Match game / Scale trainer
/ Range), the tuner meter, the steadiness readout, and the live pitch graph.}
\caption{Realtime voice monitor and vocal trainer.}
\label{fig:scr-realtime}
\end{figure}

\subsection{Transposer}

The Transposer tab shifts a score to a new key. File mode transposes an
uploaded MIDI or MusicXML file server-side and offers Camelot-compatible key
presets; hum mode transposes the last transcription client-side.

\begin{figure}[H]
\centering
\reportfig{0.85}{screenshot-transposer.png}{The Transposer tab in hum mode: the
just-hummed melody transposed client-side, with the shift control, the To-key
picker, and the transposed chords. Chord functions (the Roman numerals) are
preserved, so only the chord names move.}
\caption{Transposer: hum mode.}
\label{fig:scr-transposer}
\end{figure}

\subsection{Sing-Along and AI coaching}

The Sing-Along tab plays an uploaded score as a guide, tracks the voice live,
and scores each note. After a take it shows a deterministic analysis of the
performance offline, and, on request, sends a numeric summary to an external
language model for spoken-language coaching (one of the opt-in, non-local
features described in Section~\ref{sec:coaching}).

\begin{figure}[H]
\centering
\reportfig{0.85}{screenshot-singalong.png}{The Sing-Along tab after a take: the
per-note scoring overview, the deterministic analysis (pitch bias, drift, leap
versus step accuracy, weakest register, worst notes), and the optional
``Get coaching'' panel with its language selector.}
\caption{Sing-Along and AI coaching.}
\label{fig:scr-singalong}
\end{figure}

\subsection{Ear Trainer}

The Ear Trainer tab plays a note group and asks the user to identify it by
tapping a multiple-choice button, with sub-modes for intervals, chord qualities,
scales and modes, and scale degrees, at three difficulty levels. It is fully
client-side, with no microphone and no server.

\begin{figure}[H]
\centering
\reportfig{0.85}{screenshot-eartrainer.png}{The Ear Trainer tab: the sub-mode and
difficulty selectors, the play button, the multiple-choice answer buttons, and
the running score, streak, and per-mode accuracy history.}
\caption{Ear Trainer.}
\label{fig:scr-eartrainer}
\end{figure}

\end{document}
